% SOURCE_AUTHORITY: sga5_fr_workpass.tex, lines 9761--9801 (LF-normalized unit).
% SOURCE_SHA256: 791F4EFFC5E02832D5D77ED03518C8156D6F07E4C8238B03545DB93D883FBB28
Como la herramienta fundamental en la demostración de 4.1 es la teoría de los esquemas obtenidos por explosión, recordemos ahora algunos hechos al respecto. Denotando por
\[
f:X'\to X
\]
el $X$-esquema obtenido por la explosión de $Y$, sea
\[
\begin{tikzcd}[column sep=large,row sep=large]
Y' \arrow[r,"j"] \arrow[d,"g"'] & X' \arrow[d,"f"]\\
Y \arrow[r,"i"'] & X
\end{tikzcd}
\]
el diagrama cartesiano construido a partir de $f$ e $i$. Recordemos (cf. SGA 6 VII 3.1, cuyas notaciones adoptamos) que $j$ es una inmersión regular de codimensión $1$ definida por un ideal
\[
J'\simeq\mathcal O_{X'}(1)
\]
y que $Y'\simeq P_Y(N)$. Por otra parte, se tiene una sucesión exacta de $\mathcal O_{Y'}$-módulos
\[
0\to F\to g^*(N)\to \mathcal O_{Y'}(1)\to0.
\]

\begin{statement}{Lema 4.2}
Para todo $A_X$-módulo $M$, los homomorfismos
\[
f^*:H_Y^*(X,M)\to H_{Y'}^*(X',f^*(M))
\]
y
\[
f_*:H_{Y'}^*(X',f^*(M))\to H_Y^*(X,M)
\]
están ligados por la relación $f_*f^*=\mathrm{id}$. Además, $f_*$ es de grado $0$.
\end{statement}

\emph{Demostración.} Como el morfismo $f$ es de intersección completa relativa de dimensión relativa nula (SGA 6 VII 1.8), la última afirmación se sigue de la definición del homomorfismo de Gysin (SGA $4^{1/2}$ Cycle). Para la primera, la fórmula de proyección para el morfismo $f$
\[
f_*(x\,f^*(x'))=f_*(x)x',
\]
donde $x\in H^*(X,A)$ y $x'\in H_Y^*(X,M)$, permite reducirse a demostrar que $f_*(1)=1$. Basta para ello comprobar la afirmación análoga sobre $U=X-Y$, que es evidente, pues $f$ induce un isomorfismo de $U'=X'-Y'$ sobre $U$.

\begin{statement}{Observación 4.3}
Por supuesto, el enunciado análogo de 4.2 obtenido eliminando los soportes es también verdadero y se demuestra de la misma manera.
\end{statement}
